% Part: sets-functions-relations
% Chapter: arithmetization
% Section: cauchy

\documentclass[../../../include/open-logic-section]{subfiles}

\begin{document}

\olfileid{sfr}{arith}{cauchy}

\olsection{परिशिष्ट: कोशी क्रमिकांच्या रूपातील वास्तव संख्या}

\olref[cuts]{sec} मध्ये आपण डेडेकिंड छेद म्हणून वास्तव संख्यांची रचना
केली. या विभागात आपण एक पर्यायी रचना स्पष्ट करू. आज आपण जिला कोशी क्रमिका
म्हणतो तिची कोशी यांनी दिलेली व्याख्या या रचनेचा आधार आहे; पण त्या
व्याख्येचा उपयोग करून वास्तव संख्यांची \emph{रचना} करण्याचे श्रेय
एकोणिसाव्या शतकातील इतर लेखकांना, विशेषतः वायश्ट्रास, हाइने, मेरे आणि
कँटर यांना जाते. (या इतिहासाच्या छानशा आढाव्यासाठी
\citeauthor{OConnorRobertson:RN} \citeyear{OConnorRobertson:RN} पाहा.)

एकोणिसाव्या शतकाकडे वळण्यापूर्वी सायमन स्टेव्हिन (1548--1620) यांचा
विचार करणे उपयुक्त ठरेल. थोडक्यात, प्रत्येक वास्तव संख्येचा विचार तिच्या
दशांश विस्ताराच्या रूपात करता येतो हे स्टेव्हिन यांच्या लक्षात आले.
त्यामुळे $\sqrt{2}$ सारख्या अपरिमेय संख्येलाही एक व्यवस्थित दशांश विस्तार
असतो; त्याची सुरुवात अशी:
\[
	1.41421356237\ldots
\]
संचसिद्धान्तात दशांश विस्तारांचे प्रतिमानीकरण करणे अगदी सोपे आहे: त्यांना
$d \colon \Nat \to \Nat$ अशी फलने माना, जिथे $d(n)$ हे आपल्याला हवे
असलेले $n$ वे दशांश स्थान आहे. मग दशांश चिन्हाच्या आधी येणारा वास्तव
संख्येचा भाग (येथे फक्त $1$) हाताळण्यासाठी या मांडणीत थोडा बदल करावा
लागेल. तसेच, उदाहरणार्थ $0.999\ldots = 1$ याची खात्री करण्यासाठी आणखी
एक बदल (तुल्यता संबंध) करावा लागेल. पण संचसिद्धान्ताच्या चौकटीत
स्टेव्हिन यांच्या पद्धतीने वास्तव संख्यांची पूर्णतः काटेकोर रचना देणे
कठीण नाही.

स्टेव्हिन हा आपला मुख्य विषय नाही. (स्टेव्हिन यांच्याविषयी अधिक माहितीसाठी
\citealt{KatzKatz2012} पाहा.) पण त्याच्याशी जवळून संबंधित असा एक विचार
करू. $\sqrt{2}$ चा दशांश विस्तार थेट हाताळण्याऐवजी अधिकाधिक अचूक होत
जाणारे विस्तार घेऊन आपण $\sqrt{2}$ च्या अधिकाधिक अचूक परिमेय निकटनांची
\emph{क्रमिका} विचारात घेऊ शकतो:
\[
1, 1.4, 1.414, 1.4142, 1.41421,\ldots
\]
वास्तव संख्यांचा विचार “अधिकाधिक चांगल्या निकटनांद्वारे” करता येतो या
कल्पनेतून अंतर्दृष्टींची आणखी एक मालिका मिळते (जशी आपण $\Int$ पासून
$\Rat$ किंवा $\Nat$ पासून $\Int$ रचताना वापरली होती). त्या मूलभूत
अंतर्दृष्टी अशा:
\begin{enumerate}
	\item प्रत्येक वास्तव संख्या एका (कदाचित अनंत) दशांश विस्ताराने
	लिहिता येते.
	\item एका (कदाचित अनंत) दशांश विस्तारात सांकेतिक केलेली माहिती
	परिमेय संख्यांच्या क्रमिकेनेही सांकेतिक करता येते.
	\item परिमेय संख्यांची क्रमिका ही $\Nat$ पासून $\Rat$ कडे जाणारे
	फलन मानता येते; क्रमिकेतील $n$ वी परिमेय संख्या $f(n)$ माना.
\end{enumerate}
अर्थात, $\Nat$ पासून $\Rat$ कडे जाणारे \emph{कोणतेही} फलन आपल्याला वास्तव
संख्या देईल असे नाही. उदाहरणार्थ, हे फलन पाहा:
\[
	f(n) =\begin{cases}
			1 & \text{जर }n\text{ विषम असेल}\\
			0 &\text{जर }n\text{ सम असेल}
		\end{cases}
\]
येथील मुख्य अडचण अशी की $0,1,0,1,0,1,0,\ldots$ ही क्रमिका कोणत्याही
वास्तव संख्येकडे जवळ जाताना दिसत नाही. म्हणून एखाद्या वास्तव संख्येकडे
जवळ जाणाऱ्या क्रमिकांचाच विचार करण्यासाठी आपले लक्ष काही \emph{सीमा}
असलेल्या क्रमिकांपुरते मर्यादित केले पाहिजे.

सीमेची कल्पना \olref[his][set][limits]{sec} मध्ये आपण यापूर्वी पाहिली
आहे. पण तेथे वापरलेली \emph{अगदी तीच} व्याख्या येथे वापरता येणार नाही.
तेथील “$(\forall \epsilon>0)$” या पदावलीत वास्तव संख्यांवरील
परिमाणन गृहीत धरले होते; आणि आपण वास्तव संख्यांवरील फलनांच्या सीमा
विचारात घेत होतो. म्हणून ती व्याख्या येथे वापरणे म्हणजे वास्तव संख्या
आधीच उपलब्ध मानणे होईल; पण नेमकी त्यांचीच आपल्याला \emph{रचना} करायची
आहे. सुदैवाने, सीमेच्या अगदी जवळच्या कल्पनेने आपण काम करू शकतो.
\begin{defn}\ollabel{def:CauchySequence}
  $f: \Nat \to \Rat$ हे फलन \emph{कोशी क्रमिका} आहे तेव्हाच आणि
  तेव्हाच, जेव्हा प्रत्येक धन $\epsilon \in \Rat$ साठी $(\exists \ell
  \in \Nat)(\forall m, n > \ell)|f(m) - f(n)| < \epsilon$.
\end{defn}

सीमेची व्यापक कल्पना पूर्वीसारखीच आहे: यथार्थतेची एखादी पातळी
(जिचे मापन~$\epsilon$ ने केले आहे) हवी असेल, तर पाहण्यासाठी एक
“प्रदेश” असतो ($\ell$ पेक्षा मोठे कोणतेही आदान). आपल्या $1$, $1.4$,
$1.414$, $1.4142$, $1.41421$\ldots या क्रमिकेला सीमा आहे हेही सहज
दिसते: $\sqrt{2}$ चे निकटन $\nicefrac{1}{10^{n}}$ इतक्या त्रुटीच्या
आत हवे असेल, तर $n$ व्या पदानंतरचे कोणतेही पद पाहा.

मग वास्तव संख्या ही कोणतीही कोशी क्रमिका \emph{आहेच} असे म्हणण्याचा
विचार साहजिक वाटेल. पण $\Int$ आणि $\Rat$ यांच्या रचनांप्रमाणेच तो फार
भोळा ठरेल: कोणत्याही एका वास्तव संख्येचा निर्देश अनेक भिन्न कोशी क्रमिका
करतात. हे पाहण्याची एक सोपी पद्धत अशी. $f$ ही कोशी क्रमिका दिली असता,
$g(0)\neq f(0)$ एवढा अपवाद सोडून $g$ हे $f$ सारखेच फलन ठरवा. पहिल्या
संख्येनंतर दोन्ही क्रमिका सर्वत्र जुळत असल्यामुळे
\olref{def:CauchySequence} मध्ये वापरलेल्या अर्थाने त्यांची सीमा समान
आहे, आणि म्हणून त्या एकाच वास्तव संख्येची “व्याख्या” करतात, असे
आपल्याला शेवटी म्हणायचे असेल. त्यामुळे या कोशी क्रमिका एकच वास्तव
संख्या आहेत असे आपण खरे तर मानले पाहिजे.

म्हणून कोशी क्रमिकांवर पुन्हा एक तुल्यता संबंध परिभाषित करून वास्तव
संख्यांची ओळख त्या संबंधाखालील तुल्यतावर्गांशी करावी लागेल.
%READERNOTE{MRARITH-009: गोठवलेल्या इंग्रजी मजकुरात वास्तव संख्यांची ओळख equivalence relations शी करावी असे म्हटले आहे; पुढील वाक्ये आणि रचना स्पष्टपणे equivalence classes वापरतात. मराठी मजकुरात अपेक्षित तुल्यतावर्ग वापरले आहेत.}
सर्वप्रथम ज्याची सीमा $0$ आहे अशा फलनाची कल्पना हवी. कोणत्याही $h :
\Nat \to \Rat$ या फलनाविषयी असे म्हणा की \emph{$h$ ची सीमा $0$ आहे}
तेव्हाच आणि तेव्हाच, जेव्हा प्रत्येक धन $\epsilon \in \Rat$ साठी
$(\exists \ell \in \Nat)(\forall n > \ell)|h(n)| < \epsilon$.
\footnote{याची तुलना \olref[his][set][limits]{sec} मधील $\lim_{x
\mathord{\rightarrow}\infty}f(x) = 0$ या व्याख्येशी करा.}
शिवाय $f$ आणि $g$ ही $\Nat \to \Rat$ अशी फलने असतील, तर
$(f-g)(n) = f(n) - g(n)$ असे माना. आता अशी व्याख्या करा:
\[
	f \Realequiv g \text{ तेव्हाच आणि तेव्हाच जेव्हा $(f-g)$ ची सीमा $0$ आहे}.
\]
$\Realequiv$ हा तुल्यता संबंध आहे हे तपासावे लागेल; आणि तो तसा आहे.
मग, हवे असल्यास, $\Nat \to \Rat$ अशा सर्व कोशी क्रमिकांचे
$\Realequiv$ खालील तुल्यतावर्ग म्हणजे वास्तव संख्या, अशी व्याख्या करता
येईल.

\begin{prob}
प्रत्येक $n$ साठी $f(n) = 0$ असे माना. $g(n) =
\frac{1}{(n+1)^2}$ असे माना. ही दोन्ही फलने कोशी क्रमिका आहेत, आणि
दोन्ही फलनांची सीमा $0$ आहे, हे दाखवा; त्यामुळे $f \sim_\Real g$
हेही दाखवा.
\end{prob}

हे केल्यानंतर, नेहमीप्रमाणे, $f$ हा !!a{element} असलेल्या तुल्यतावर्गासाठी
आपण $\equivrep{f}{\Realequiv}$ असे लिहू. मात्र मजकूर वाचनीय ठेवण्यासाठी
पुढे अधोलिखित भाग वगळून फक्त $\equivrep{f}{}$ असे लिहू. प्रत्येक $q
\in \Rat$ साठी $q_{\Real} = \equivrep{c_{q}}{}$ असेही आपण ठरवतो,
जिथे प्रत्येक $n \in \Nat$ साठी $c_q(n) = q$ असे स्थिर फलन $c_{q}$
आहे. मग वास्तव संख्यांवरील मूलभूत संबंध आणि क्रिया परिभाषित करतो;
उदाहरणार्थ:
\begin{align*}
	\equivrep{f}{} + \equivrep{g}{} &= 	\equivrep{(f + g)}{} \\
	\equivrep{f}{} \times 	\equivrep{g}{} &= \equivrep{(f \times g)}{}
\end{align*}
जिथे $(f + g)(n) = f(n) + g(n)$ आणि $(f \times g)(n) = f(n) \times
g(n)$. $f$ आणि $g$ कोशी क्रमिका असतील तेव्हा $(f + g)$, $(f-g)$ आणि
$(f\times g)$ यांपैकी प्रत्येकही कोशी क्रमिका आहे हे अर्थात तपासावे
लागेल; तसे ते आहेत, आणि ही सिद्धता तुमच्यासाठी सोडली आहे.
%READERNOTE{MRARITH-012: गोठवलेला मजकूर येथे बेरीज, वजाबाकी आणि गुणाकाराने कोशी क्रमिका मिळतात एवढीच संवृति तपासायला सांगतो. तुल्यतावर्गांवरील क्रिया सुव्याख्यायित होण्यासाठी प्रतिनिधी बदलल्यावर वर्ग बदलत नाही हेही तपासणे आवश्यक आहे; तो स्वतंत्र तपशील स्रोतामध्ये दिलेला नाही.}

शेवटी क्रमाची संकल्पना परिभाषित करू. $\equivrep{f}{}$ हा वर्ग
\emph{धन} आहे तेव्हाच आणि तेव्हाच, जेव्हा $\equivrep{f}{}\neq 0_\Real$
आणि $(\exists \ell \in \Nat)(\forall n > \ell)0 < f(n)$
या दोन्ही अटी पूर्ण होतात. मग $\equivrep{f}{} <
\equivrep{g}{}$ तेव्हाच आणि तेव्हाच, जेव्हा $\equivrep{(g - f)}{}$
धन आहे, असे म्हणा.
%READERNOTE{MRARITH-010: गोठवलेल्या इंग्रजीतील धनतेच्या व्याख्येत वास्तव तुल्यतावर्गाची तुलना 0_Rat शी केली आहे. येथे त्याच रचनेतील वास्तव शून्य 0_Real वापरले आहे.}
ही व्याख्या सुव्याख्यायित आहे, म्हणजे ती तुल्यतावर्गातून निवडलेल्या
“प्रतिनिधी” फलनावर अवलंबून नाही, हे तपासावे लागेल.
%\begin{align*}
%	\equivrep{f}{} \leq \equivrep{g}{} \text{ तेव्हाच आणि तेव्हाच जेव्हा }&\equivrep{f}{} = \equivrep{g}{} \text{ किंवा }\\
%	&(\exists \ell \in \Nat)(\forall n \in \Nat)(\ell < n \lif f(n) \leq g(n))
%\end{align*}
हे केल्यानंतर या व्याख्यांमुळे योग्य बैजिक गुणधर्म मिळतात हे दाखवणे
बरेच सोपे आहे; म्हणजे:
\begin{thm}\ollabel{thm:cauchyorderedfield}
कोशी क्रमिकांचे वरील तुल्यतावर्ग एक क्रमित क्षेत्र बनवतात.
\end{thm}
%READERNOTE{MRARITH-011: गोठवलेल्या इंग्रजीतील प्रमेय, पुढील प्रश्न आणि संपूर्णतेचे विधान Cauchy sequences हेच क्षेत्र बनवतात अशी संक्षिप्त भाषा वापरतात. प्रत्यक्ष क्षेत्राचे घटक Realequiv संबंधाखालील तुल्यतावर्ग आहेत; मराठी मजकुरात ते स्पष्ट केले आहे. पुढील सिद्धतेत जेव्हा क्रमिकेवर वास्तव-क्रम लावला जातो तेव्हा तिचा तुल्यतावर्ग अभिप्रेत आहे.}

\begin{proof}
सराव.
\end{proof}

\begin{prob}
कोशी क्रमिकांचे तुल्यतावर्ग एक क्रमित क्षेत्र बनवतात हे सिद्ध करा.
\end{prob}

अशा रीतीने रचलेल्या वास्तव संख्यांना संपूर्णता गुणधर्म आहे हे सिद्ध
करणे अधिक कठीण आहे; म्हणून त्याची सिद्धता आपण देऊ.

\begin{thm}
ऊर्ध्वबंध असलेल्या कोशी क्रमिकांच्या तुल्यतावर्गांच्या प्रत्येक अरिक्त
संचाला लघुतम ऊर्ध्वबंध असतो.
\end{thm}

\begin{proof}[सिद्धतेचा आराखडा]
$S$ हा ऊर्ध्वबंध असलेल्या कोशी क्रमिकांचा कोणताही अरिक्त संच असू द्या.
म्हणून काही $p \in \Rat$ असे आहे की $p_{\Real}$ हा $S$ चा ऊर्ध्वबंध
आहे. $r \in S$ असू द्या; मग काही $q \in \Rat$ असे आहे की
$q_{\Real} < r$. त्यामुळे $S$ ला लघुतम ऊर्ध्वबंध असेल, तर तो
$q_\Real$ आणि $p_\Real$ यांच्या दरम्यान (दोन्ही टोकांसह) असेल.

लघुतम ऊर्ध्वबंधाकडे खालून आणि वरून एकाच वेळी जवळ जात आपण तो शोधू.
विशेषतः, $f, g \colon \Nat \to \Rat$ अशी दोन फलने परिभाषित करू:
$f$ ने वरून आणि $g$ ने खालून लघुतम ऊर्ध्वबंधाकडे जवळ जावे हा उद्देश
आहे. सुरुवात अशा व्याख्यांनी करू:
\begin{align*}
	f(0) &= p \\
	g(0) &= q
\end{align*}
मग $a_n = \frac{f(n) + g(n)}{2}$ असेल, तर पुढील व्याख्या
करा:\footnote{ही पुनरावर्ती व्याख्या आहे. पण पुनरावर्ती व्याख्या ग्राह्य
आहेत असे मानण्याचे कोणतेही कारण आपण \emph{अजून} दिलेले नाही.}
\begin{align*}
	f(n+1) &=
	\begin{cases}
		a_n &\text{जर }(\forall h \in S)\equivrep{h}{} \leq (a_n)_\Real\\
		f(n)&\text{अन्यथा}
	\end{cases}\\
	g(n+1) &=
	\begin{cases}
		a_n &\text{जर }(\exists h \in S)\equivrep{h}{} \geq (a_n)_\Real\\
	 	g(n) &\text{अन्यथा}
	\end{cases}
\end{align*}
$f$ आणि $g$ या दोन्ही कोशी क्रमिका आहेत. (हे बऱ्यापैकी सहज तपासता
येते, पण तो सराव तुमच्यासाठी सोडला आहे.) प्रत्येक पायरीवर $f$ आणि
$g$ मधील फरक निम्मा होत असल्यामुळे $(f-g)$ या फलनाची सीमा $0$ आहे.
त्यामुळे $\equivrep{f}{} = \equivrep{g}{}$.

सर्वप्रथम $\equivrep{f}{}$ हा $S$ चा ऊर्ध्वबंध आहे, म्हणजे
$(\forall h \in S)\equivrep{h}{} \leq \equivrep{f}{}$, हे दाखवू.
(वाटेत आपण \olref{thm:cauchyorderedfield} वापरू.) $h \in S$ असू द्या
आणि व्याघातासाठी $\equivrep{f}{} < \equivrep{h}{}$ असे समजा; म्हणजे
$0_\Real < \equivrep{(h-f)}{}$. $f$ ही एकस्वनिक घटती कोशी क्रमिका
असल्यामुळे काही $n \in \Nat$ असे आहे की
$\equivrep{(c_{f(n)} - f)}{} < \equivrep{(h-f)}{}$. म्हणून:
\[
	(f(n))_\Real = \equivrep{c_{f(n)}}{} < \equivrep{f}{} + \equivrep{(h-f)}{} = \equivrep{h}{},
\]
पण रचनेनुसार $\equivrep{h}{} \leq (f(n))_\Real$ असल्यामुळे हा
व्याघात आहे.

आता $\equivrep{f}{} = \equivrep{g}{}$ हाच $S$ चा \emph{लघुतम}
ऊर्ध्वबंध आहे हे दाखवू. $j$ ही कोणतीही कोशी क्रमिका असू द्या आणि
$\equivrep{j}{} < \equivrep{g}{}$ असे समजा. वरच्याप्रमाणेच युक्तिवाद
केल्यास ($g$ हे \emph{वाढते} आहे या तथ्याचा उपयोग करून) काही $n \in
\Nat$ असे आहे की $\equivrep{j}{} < (g(n))_\Real$. पण रचनेनुसार काही
$h \in S$ असे आहे की $(g(n))_\Real \leq \equivrep{h}{}$; म्हणून
$\equivrep{j}{} < \equivrep{h}{}$. त्यामुळे $\equivrep{j}{}$ हा
$S$ चा ऊर्ध्वबंध नाही.
\end{proof}

\end{document}
